\subsection{Two distinct empirical objects}

The literature commonly described as ``LLM alpha'' contains two different
empirical objects.  Formulaic and cross-sectional systems generate, search,
score, or select an executable factor or factor portfolio.  We denote these
systems by stratum \(F\).  Their output can usually be represented as a dated
cross-sectional score,
\[
    s_t = g_\theta(X_{\leq t},\mathcal{O},\mathcal{M}_t)
    \in \mathbb{R}^{N_t},
\]
where \(X_{\leq t}\) is the information available by formation date \(t\),
\(\mathcal{O}\) is an operator or strategy language, and \(\mathcal{M}_t\)
denotes any search memory or feedback state.  Alpha-GPT, FAMA, AlphaAgent,
Alpha-Jungle, Chain-of-Alpha, FactorMiner, Hubble, and AlphaCrafter illustrate
the progression from prompted formula generation to constrained program
search, memory, and adaptive factor-to-execution pipelines
\citep{li2023alphagpt,LiEtAl2024FAMA,TangEtAl2025AlphaAgent,shi2025alphajungle,cao2025chainofalpha,wang2026factorminer,shi2026hubble,yuan2026alphacrafter}.
Recent work also combines LLM screening with conventional portfolio
estimation, making the stock-selection step, rather than only the weight
optimizer, an explicit object of study \citep{caner2026agenticscreening}.

Sequential trading and portfolio systems, denoted by stratum \(T\), instead
emit dated orders, positions, or portfolio weights.  Their empirical object is
a policy,
\[
    (a_t,w_t) = \pi_\theta(h_t,\mathcal{M}_t,\mathcal{T}_t),
\]
where \(h_t\) is the observable history, \(\mathcal{M}_t\) is an evolving
memory or reflection state, and \(\mathcal{T}_t\) is the set of tools available
at date \(t\).  FinMem and FinAgent emphasize memory, multimodal perception,
and tool use; FinCon and TradingAgents add manager--analyst hierarchies or
role-specialized deliberation; and more recent systems add prompt adaptation,
position-aware risk control, or direct portfolio construction
\citep{YuEtAl2024FinMem,zhang2024finagent,YuEtAl2024FinCon,XiaoEtAl2024TradingAgents,papadakis2025atlas,tian2025tradinggroup,liu2025finpos,guo2025mass}.
The distinction is operational rather than semantic.  An \(F\) system can in
principle be placed on a common point-in-time panel and operator contract.  A
\(T\) system must be replayed through its entire observe--reason--act--execute
path because its next decision depends on earlier actions and state.  Replacing
either object with a thematically similar published characteristic does not
reproduce the generator, search process, memory, execution rule, or dated
output.  We therefore call such substitutions \emph{mechanism-inspired
proxies}, not replications.

This distinction complements the asset-pricing replication literature.
Published predictor returns tend to attenuate out of sample and after
publication, while large-scale reassessments show that specification,
construction, and multiple-testing choices materially affect which anomalies
survive \citep{McLeanPontiff2016,HarveyLiuZhu2016,hou2020replicating,ChenZimmermann2022,JensenKellyPedersen2023}.
Agentic systems add further researcher degrees of freedom---model version,
prompt, tools, memory, search budget, stochastic seed, and execution
semantics---to the familiar data-mining problem.  White's reality check,
Hansen's superior-predictive-ability test, and the deflated Sharpe ratio thus
remain relevant even when the search procedure is expressed in natural
language rather than in a conventional optimizer
\citep{White2000RealityCheck,Hansen2005SPA}.

\subsection{Frozen census protocol}

We constructed a public-system census through 2026-08-02 23:59 UTC.  The unit
of observation is a \emph{named system/version lineage}, not a paper, a
repository, a benchmark row, or a citation.  Publications, repositories,
model cards, and demonstrations are linked many-to-one to a lineage.  Searches
covered arXiv categories cs.AI, cs.CL, cs.LG, cs.CE, q-fin.CP, q-fin.PM,
q-fin.ST, and q-fin.TR; the ACL Anthology; OpenReview; the ACM Digital Library;
and SSRN.  Crossref and OpenAlex records were used to reconcile titles and
publication versions.  GitHub was used only for exact-title, author-linked, or
publication-linked artifacts.  We combined three prespecified query families:
LLM/agent terms with alpha mining, factor discovery, formulaic alpha, or
strategy generation; LLM/agent terms with trading, portfolio, or stock
selection and a return, Sharpe, backtest, or live-evaluation term; and
benchmark, audit, reproducibility, or leakage terms with alpha mining or
trading agents.  Backward and forward searches began from AlphaBench,
AlphaQT-Bench, FINSABER, and the Agentic Trading audit
\citep{alphabench2026,LuoEtAl2026AlphaQT,LiEtAl2026FINSABER,XiaEtAl2026AgenticTrading}.
Curated lists were discovery aids only and never counted as systems.

Inclusion required an operational investment output.  An \(F\) lineage had to
generate, score, search, or select an executable factor, cross-sectional score,
or factor portfolio.  A \(T\) lineage had to emit dated tradable actions or
weights and report a closed-loop historical, paper-trading, or live
evaluation.  Static question answering, report generation, and forecasts with
no tradable output were excluded.  Generic financial assistants and private
claims inaccessible by the cutoff were also excluded.  Endogenous market
simulations were retained as evaluation environments but not as candidate
alpha systems.  Exclusions and boundary decisions remain in a separate ledger
so that the denominator cannot be revised after observing replication outcomes.

The registry separates five mutually exclusive strata.  Strata \(F\) and
\(T\) form the planned system denominator.  Stratum \(B\) contains benchmarks,
audits, and evaluation environments; \(C\) contains public community software
without a sufficiently specified scholarly system claim; and \(M\) contains
non-LLM or mechanistic comparators.  Table~\ref{tab:census-strata} reports the
frozen counts.  The resulting corpus has 103 lineages, of which 67 are in the
primary \(F+T\) denominator.  The remaining 36 records contextualize artifact
availability and evaluation but are not treated as candidate agent returns.
The earlier 55-reference working list mixed papers, repositories, benchmarks,
and systems and therefore was not used as a denominator.  Reconciliation and
systematic searching added 48 named lineages to that seed material; a later
MarketSenseAI validation paper was linked to its existing lineage rather than
counted again.

Here ``frozen'' refers to the census cutoff, lineage dispositions, and fixed
denominators.  It does not imply that the later return analysis remained an
untouched holdout; the post-outcome evaluator repairs and limited-liability
amendment are disclosed in the common-task design.

\begin{table}[t]
    \centering
    \small
    \caption{Frozen public-system census at the 2026-08-02 cutoff.}
    \label{tab:census-strata}
    \begin{tabular}{clrc}
        \toprule
        Stratum & Operational definition & \(N\) & In main set \\
        \midrule
        \(F\) & Formulaic/cross-sectional alpha & 29 & Yes \\
        \(T\) & Sequential trading/portfolio & 38 & Yes \\
        \(B\) & Benchmark/audit/environment & 23 & No \\
        \(C\) & Community software & 5 & No \\
        \(M\) & Mechanistic comparator & 8 & No \\
        \midrule
        All & Named system/version lineages & 103 & 67 \\
        \bottomrule
    \end{tabular}
\end{table}

Deduplication was substantive rather than string-based.  The Alpha-GPT arXiv
preprint and EMNLP demonstration are two publications of its first release;
Alpha-GPT~2.0 is recorded as a successor release within the same lineage
\citep{li2023alphagpt,WangEtAl2025AlphaGPT,YuanEtAl2024AlphaGPT2}.  FINSABER's main and
reproduction branches are versioned artifacts of one benchmark, and the arXiv
and announced WWW manifestations for Agent Market Arena are one benchmark lineage
\citep{LiEtAl2026FINSABER,zhang2026agentmarketarena}.  The later MarketSenseAI
validation study is evidence attached to the existing MarketSenseAI lineage,
not a new system \citep{FatourosMetaxas2026SignalOrNoise}.  Conversely, FactFin and
FinLake-Bench (also called FinLeak-Bench in the paper body) are separate named
operational objects---a strategy framework and a leakage benchmark---even
though they share one publication \citep{li2025profitmirage}.  AlphaForge and
AlphaGen remain distinct methods, and the two QuantAgent records remain
distinct because they have different authors, tasks, and architectures.
Finally, Agent Market Arena, Agent Trading Arena, and TradeArena are distinct:
respectively a live multi-market benchmark, an endogenous zero-sum market
simulation, and an auditable exogenous-market testbed
\citep{zhang2026agentmarketarena,ma2025agenttradingarena,xue2026tradearena}.

Across the frozen registry, artifact audit, and native-fidelity ledger, we
record every lineage's primary publication, official-artifact basis, cutoff
revision, license, native task, output contract, evaluator, execution attempt
where applicable, and fidelity.
Artifact availability is graded from no public artifact path (R0), through a
descriptive artifact (R1) and executable components (R2), to a packaged
runnable structure (R3).  These are static packaging tiers, not evidence that
an artifact executed successfully or produced dated returns.  Fidelity is classified separately as
official-native, official-adapted, independent reimplementation,
mechanism-inspired, unavailable, legally unusable, or technically
incompatible.  A failed installation is therefore an artifact failure, not a
zero return; a paper-only record is not evidence that the underlying method
lacks alpha; and a mechanism-inspired proxy is not promoted to an agent
replication.

\subsection{What external evaluations establish}

The benchmark literature does not support treating reported headline returns
as commensurable estimates.  AlphaQT-Bench separates syntax and executability
from causal safety, functional correctness, and structural compliance,
showing why a formula that parses is not necessarily a valid quantitative
program \citep{LuoEtAl2026AlphaQT}.  AlphaEval likewise evaluates formula
alphas through a dedicated execution and financial-quality pipeline rather
than accepting surface-form plausibility \citep{ding2025alphaeval}.  These
benchmarks motivate a first-stage validity gate before any economic comparison.

Longer and more standardized trading evaluations generally report substantial
attenuation.  FINSABER evaluates LLM investment strategies over more than 20
years and over 100 symbols and reports that advantages observed in short
experiments deteriorate over the longer horizon \citep{LiEtAl2026FINSABER}.
StockBench evaluates 82 trading days from March through June 2025 with three
runs per configuration and reports that most evaluated agents fail to beat
buy-and-hold \citep{ChenEtAl2025StockBench}.  KTD-Fin evaluates ten agents on the
CSI~300 over 2024--2026; after masking asset identities and applying Barra-style
attribution, it reports that much apparent return reflects passive exposure or
style tilts rather than stock-selection alpha \citep{ZhuEtAl2026KTDFin}.  PortBench
spans six asset classes and ten years and reports that 90\% of evaluated
model--profile combinations fail an equal-weight benchmark, with severe losses
in stress episodes \citep{ZhaoChenSu2026PortBench}.  These results establish that
evaluation horizon, baseline, identity masking, factor attribution, and stress
coverage can reverse a ranking.  They do not establish that every omitted
system has zero alpha, because the benchmarks differ in assets, prompts,
available tools, model versions, and execution contracts.

Artifact audits identify a separate evidentiary bottleneck.  The Agentic
Trading audit screened 92 candidate papers and retained 77, but classified only
19 as closed-loop trading studies.  Among those 19, it found two with temporal
splits, one with transaction costs, one addressing universe or survivorship
construction, and 11 specifying execution semantics; 15 were paper-only and
none reached its exact native-reproduction tier \citep{XiaEtAl2026AgenticTrading}.
Those counts describe reporting and artifact availability under the audit's
scope; they are not performance estimates.  They nevertheless show why a
paper-count denominator and a runnable-system denominator answer different
questions.

Profit Mirage provides direct evidence on temporal contamination, but its
causal interpretation requires care.  Under the authors' implementation, five
well-known agents re-evaluated with GPT-4o exhibit Sharpe-ratio declines of
51.48\%--62.23\% and total-return declines of 50.18\%--71.85\% after the
model's stated knowledge cutoff.  In counterfactual tests, 69.03\%--82.13\% of
actions remain unchanged after material input perturbations, and the companion
2,000-item FinLake/FinLeak benchmark finds high recall of historical market
facts \citep{li2025profitmirage}.  These findings are consistent with
memorization contributing to backtest performance.  They do not by themselves
identify the causal share of decay attributable to pretraining leakage rather
than regime change, implementation differences, or model--tool interactions.
Our design therefore treats the model cutoff as a contamination risk and uses
point-in-time inputs under the disclosed amended evaluator; it does not subtract a presumed ``leakage effect''
from returns.

The evidence is not uniformly negative.  A longitudinal MarketSenseAI study
reports live-generated S\&P~500 signals over 19 months with average monthly
return of 2.18\% versus 1.15\% for the benchmark, a Monte Carlo
\(p\)-value of 0.003, and an information-coefficient information ratio of
0.489 with \(p=0.024\); its 35-month S\&P~100 comparison is positive but not
statistically significant \citep{FatourosMetaxas2026SignalOrNoise}.  Because this is a
follow-up from the system's own research lineage rather than an independent
replication, we treat it as external-time evidence, not independent
verification.  In a more constrained comparison, Macro Economists in the
Machine holds the macro inputs and portfolio engine fixed across LLM and rule
agents over 124 weekly rebalances.  It reports Sharpe improvements of 0.044 and
0.040 for the Hawkish and Debate agents relative to the rule agent (both block
bootstrap \(p<0.10\)) and a net advantage surviving one-way trading costs up
to 30 basis points.  The Debate agent does not beat the best single agent
(difference \(-0.004\), \(p=0.769\)); the sample covers one rate cycle and the
reported \(p\)-values are not adjusted for multiple comparisons
\citep{wang2026macroeconomists}.  Fin-Analyst supplies a further live common-
window observation: over the short FinMMEval evaluation it reports a 13.51\%
TSLA return, 28.33 percentage points above buy-and-hold, while its BTC system
finishes approximately flat and the asset ranking reverses relative to the
interim leaderboard \citep{rashid2026finanalyst}.  The reversal is evidence
about short-window ranking instability as much as it is evidence about the
agent.

Taken together, the literature supports an artifact-to-evidence attrition
design rather than a league table of incomparable Sharpe ratios.  We audit
every one of the 67 \(F+T\) lineages under the frozen rules, preserve failures,
attempt native execution when the public requirements make that attempt
identified and proportionate, report native-protocol evidence when it is
recoverable, and restrict common-task comparisons to outputs that satisfy the
same point-in-time information, execution, and cost contract.  This separates
three questions that prior work often conflates: whether a public artifact
exists, whether it can generate the claimed dated output, and whether that
output earns robust net alpha under a common evaluator.
